% Independent mathematical audit notes.  No preamble; all assertions below
% concern linear cones unless a separate geometric hypothesis is stated.
% Standard foundations: Farkas alternatives, rational LP, basic feasible
% solutions, rank-nullity, and elementary face theory.  The application and
% the explicit transfer/family statements should not be advertised as a new
% discovery of strict complementarity or polynomial-time linear programming.

\section{Complete support by one primal--dual feasibility problem}

Let $A\in\mathbb Q^{m\times k}$ and write
\[
 C(A)=\{x\in\mathbb R^k:Ax=0,\ x\geq0\},\qquad
 S(A)=\{j:\text{some }x\in C(A)\text{ has }x_j>0\}.
\]
The complement $F(A)=\{1,\ldots,k\}\setminus S(A)$ is the set of
coordinates forced to zero.  All inequalities between vectors are
coordinatewise.  A dual exposing vector means a vector
$y=A^T\lambda\geq0$ whose zero set is exactly $S(A)$.

\begin{theorem}[Complete support from a joint linear certificate]
The linear system
\begin{equation}\label{eq:audit-joint}
 Ax=0,\qquad y=A^T\lambda,\qquad x\geq0,\quad y\geq0,
 \qquad x+y\geq\mathbf1
\end{equation}
is feasible for every rational matrix $A$.  Every feasible solution has
\[
 \operatorname{supp}(x)=S(A),\qquad
 \operatorname{supp}(y)=F(A).
\]
In particular, a verifier needs only exact matrix arithmetic and sign
tests, rather than access to the optimization routine producing the
certificate.
\end{theorem}

\begin{proof}
First prove the coordinate alternative.  Let $a_j$ be column $j$ of
$A$.  There is an element of $C(A)$ with positive $j$th coordinate if
and only if
\[
 A_{-j}u=-a_j,\qquad u\geq0
\]
is feasible: scale the positive coordinate to one.  If this system is
infeasible, Farkas' lemma supplies $\lambda^{(j)}$ such that
\[
 A_{-j}^T\lambda^{(j)}\geq0,
 \qquad a_j^T\lambda^{(j)}>0.
\]
Thus $A^T\lambda^{(j)}\geq0$, with a positive $j$th coordinate.
Conversely, such a vector forces $x_j=0$ for every $x\in C(A)$, since
\[
 0=(Ax)^T\lambda^{(j)}=x^TA^T\lambda^{(j)}
\]
is a sum of nonnegative terms.  This also checks the sign in the Farkas
alternative.

For each $j\in S(A)$ choose a primal witness and scale it so its $j$th
coordinate is at least one.  Their sum $x$ is positive on every
coordinate of $S(A)$ and zero outside $S(A)$.  For each $j\in F(A)$
choose the dual witness just constructed and scale it so its $j$th
coordinate is at least one.  Sum the corresponding multipliers to obtain
$\lambda$, and set $y=A^T\lambda$.  Every summand of $y$ vanishes on
$S(A)$ by orthogonality with feasible primal witnesses.  Consequently
$x+y\geq\mathbf1$, proving existence.  Empty sums cause no difficulty.

Now consider any solution of \eqref{eq:audit-joint}.  We have
$x^Ty=0$, so nonnegativity gives $x_jy_j=0$ for each $j$.
The last inequality ensures that exactly one of $x_j,y_j$ is positive.
Every positive $x_j$ witnesses $j\in S(A)$, while every positive $y_j$
forces $j\notin S(A)$.  This proves the claimed partition.
\end{proof}

The theorem is a specialization of classical alternatives and strict
complementarity for the nonnegative orthant.  Its use here is a complete
replacement for sufficient one-row or two-row sign propagation.

\subsection{Polynomial bit complexity, including redundant rows and lineality}

\begin{proposition}[Small exact certificates]
After clearing row denominators, suppose $A$ is integral with
$|A_{ij}|\leq H$, where $H\geq1$.  System
\eqref{eq:audit-joint} has a rational solution of polynomial encoding
length in $m,k,\log H$.  More explicitly, put $v=3k+2m$ and let
$\rho\leq m+2k$ be the rank of the equality matrix in the standard-form
system below.  One may choose a common denominator and all primal and
dual variable numerators of magnitude at most
\[
 D=\rho!H^\rho,
\]
except that a multiplier $\lambda_i$ obtained as a difference of two
nonnegative variables may require the bound $2D$.  The convention
$D=1$ applies if $\rho=0$.  In particular the total number of bits is
$O(v\rho(\log(H+1)+\log(\rho+1)))$, up to ordinary bookkeeping.
\end{proposition}

\begin{proof}
Split $\lambda=\lambda^+-\lambda^-$ and introduce a nonnegative surplus
vector $z$.  All variables in
\[
 Ax=0,\qquad A^T\lambda^+-A^T\lambda^- - y=0,
 \qquad x+y-z=\mathbf1,
 \qquad x,y,z,\lambda^+,\lambda^-\geq0
\]
are nonnegative.  This is a feasible standard-form rational system
$Mu=b$, with integral coefficients bounded by $H$ and with entries of
$b$ equal to zero or one.  Remove redundant equality rows, leaving
$\rho$ independent rows.  Row removal means choosing a subset of the
original rows; it does not enlarge the coefficient bound.  Since the
system is feasible, the removed equations follow from those retained
with the correct right-hand sides.

For completeness, a nonempty standard-form feasible set has a basic
feasible solution even when the original free-multiplier formulation
has lineality.  Choose a feasible vector with minimum positive support.
If the columns of $M$ on that support were dependent, perturb within
their dependence in a direction and for a distance that makes one
positive coordinate zero while keeping all others nonnegative.  This
preserves the equalities and contradicts minimality.  The support
columns are therefore independent.  Extend them to a basis of $\rho$
columns, assigning zero to any added coordinates.  This produces a
basic feasible solution.  Equivalently, the nonnegative standard-form
feasible set contains no affine line, despite potentially unbounded
recession directions.

Cramer's rule on the corresponding $\rho\times\rho$ nonsingular
integer matrix bounds its determinant and all numerator determinants
by $\rho!H^\rho$.  The determinant is a nonzero integer.  This proves
the bounds and the claimed encoding length.  Recombining the split
multipliers introduces at most the stated factor two.

For rational $A$, clearing denominators row by row has polynomial bit
cost in its input encoding.  If $A'=RA$ is the resulting integer matrix,
where $R$ is diagonal, a multiplier $\mu$ for $A'$ corresponds to the
source multiplier $R\mu$ for $A$.  This transformation again has
polynomial encoding length.
\end{proof}

Classical rational linear-programming algorithms therefore find such a
certificate in polynomial bit time.  This statement concerns the
mathematical LP primitive; it does not assert a polynomial pivot bound
for a chosen simplex implementation.  Clearing a common positive
denominator in a joint certificate also produces integral $x,y,\lambda$:
the inequality $x+y\geq\mathbf1$ becomes stronger after this scaling.

\section{The support restriction is the intrinsic linear span}

\begin{proposition}[Effective dimension]
Let $S=S(A)$.  Extension by zero gives a cone isomorphism
\[
 C(A_S)\longrightarrow C(A),
\]
and
\[
 \operatorname{span} C(A_S)=\ker A_S,
 \qquad \dim\operatorname{span} C(A)=|S|-\operatorname{rank} A_S.
\]
\end{proposition}

\begin{proof}
All deleted coordinates vanish on $C(A)$, so restriction and extension
by zero are inverse linear maps on the cones.  There is a vector
$p\in\ker A_S$ with $p>0$.  Given $u\in\ker A_S$, choose
$\epsilon>0$ small enough that $p\pm\epsilon u\geq0$.  Both of these
vectors belong to $C(A_S)$, and their difference is $2\epsilon u$.
This proves the span assertion.  Rank-nullity gives the formula.
\end{proof}

Extension by zero also preserves the complete face lattice, extreme
rays, linear objectives and the integral lattice in these coordinate
spaces.  More generally it preserves any further conditions on the
retained coordinates.  When applying this to normal surfaces, one must
still justify the relationship between the matching cone used for
support computation and the actual standard-coordinate source cone.
Forced zeros in a relaxation safely remain forced in every geometric
subcone; maximality in a relaxation does not establish maximality after
quadrilateral exclusions or other geometric restrictions.

\section{Exact reuse when allowed columns are deleted or added}

Throughout this section, all columns belong to one fixed matrix of
source constraints.  The results do not authorize reuse across a
Pachner move or any other change to the source equations without a
separate transport proof.  Write $C_I=C(A_I)$ for an allowed column set
$I$.  Let $S=S(A_I)$, $F=I\setminus S$, and choose a complete exposing
multiplier $\lambda$ such that
\[
 y_i=a_i^T\lambda=0\ (i\in S),\qquad
 y_i>0\ (i\in F).
\]

\subsection{Deletion and lifting a smaller dual certificate}

\begin{proposition}[Restriction and dual lifting]
For $I'\subseteq I$, every coordinate of $F\cap I'$ remains forced
zero.  Complete support for $I'$ may therefore be computed using only
$R=S\cap I'$.  If $I'$ contains $S$, the original support and effective
dimension are unchanged.  In general the effective dimension cannot
increase.

Suppose $\eta$ is a complete exposing multiplier for the smaller cone
$C(A_R)$, and put $v_i=a_i^T\eta$ for all $i\in I'$.  Define
\[
 \alpha=1+\max_{i\in F\cap I'}
              \frac{\max(0,-v_i)}{y_i},
\]
where the maximum of an empty set is zero.  Then
$\eta+\alpha\lambda$ is a complete exposing multiplier on $I'$.
All these operations preserve polynomial certificate bit length.
\end{proposition}

\begin{proof}
The restriction of $y$ is still a nonnegative row-space consequence,
so its positive coordinates remain zero on every feasible vector.
This gives the cone isomorphism with $C(A_R)$.  The cone for $I'$ is a
coordinate face of the old cone, embedded with zeros, so its dimension
cannot increase.  If only old forced coordinates are deleted, the cones
are identical under that embedding.

On $R$, the old exposing vector is zero, so the new exposing vector
equals $v$.  On $F\cap I'$, its coefficient is
$v_i+\alpha y_i\geq y_i>0$.  Its zero set is therefore precisely the
new maximal support.  Rational maxima, products and quotients of the
displayed certificate entries have polynomial bit length.
\end{proof}

\subsection{Which added columns can invalidate an exposing vector?}

\begin{theorem}[Dual sign test for extensions]
Add columns indexed by $E$, disjoint from $I$, and let
\[
 w_e=a_e^T\lambda,\qquad
 E_+=\{e:w_e>0\},\quad E_0=\{e:w_e=0\},\quad
 E_-=\{e:w_e<0\}.
\]
Every $q\in C_{I\cup E}$ satisfies the exact balance identity
\begin{equation}\label{eq:audit-portal}
 \sum_{i\in F}y_iq_i+\sum_{e\in E_+}w_eq_e
       =\sum_{e\in E_-}(-w_e)q_e.
\end{equation}
Consequently, positive use of an old forced coordinate or of a column
in $E_+$ requires positive use of at least one column in $E_-$.

If $E_-=\varnothing$, every coordinate in $F\cup E_+$ is still forced
zero, and the expanded cone is isomorphic to $C(A_{S\cup E_0})$.
Writing $d=\dim\operatorname{span}C_I$ and $d'$ for the corresponding
dimension after extension, one has
\[
 d\leq d'\leq d+|E_0|.
\]
If every added column has strictly positive $w_e$, the old cone and
support remain unchanged.
\end{theorem}

\begin{proof}
Apply $\lambda^T$ to the matching equation for $q$ and move its negative
terms to the other side.  This proves \eqref{eq:audit-portal}.
Every summand on both sides is nonnegative, giving the necessity
statement and, when $E_-$ is empty, all asserted forced zeros.
Old feasible vectors extend by zero, so $d'\geq d$.  Since
$\dim\ker A_S=d$, adding $|E_0|$ columns increases the kernel dimension
by at most $|E_0|$.  The new feasible span lies in that kernel, yielding
the upper bound.  The last assertion follows when $E_0$ is empty.
\end{proof}

\begin{example}[One added column can cause arbitrary revival]
Let $A_I=(1,\ldots,1)$ have $r$ columns.  Its nonnegative kernel is
zero.  Add the single column $-1$.  The new cone is
\[
 \{(u_1,\ldots,u_r,t)\geq0:t=u_1+\cdots+u_r\},
\]
which has all $r+1$ coordinates in its maximal support and dimension
$r$.  The old multiplier $\lambda=1$ evaluates negatively on the added
column.  Thus no dimension bound depending only on the old effective
dimension and the number of added columns can hold without additional
hypotheses.  In particular the nonnegative sign hypothesis above is
substantive.
\end{example}

\subsection{Contracting the live core}

\begin{theorem}[Quotient extension and exact dimension sequence]
Let $S$ be a column set admitting $p>0$ with $A_Sp=0$.  Set
$V=\operatorname{im}A_S$, and choose an exact linear map $Q$ with
$\ker Q=V$.  For any allowed set $J\supseteq S$, put
\[
 R=J\setminus S,\qquad B=QA_R.
\]
The coordinate projection onto $R$ maps $C(A_J)$ surjectively onto
$C(B)$.  If $U=S(B)$, then
\[
 S(A_J)=S\cup U.
\]
There is a short exact sequence of real vector spaces
\[
 0\longrightarrow\ker A_S
   \xrightarrow{\;u\mapsto(u,0)\;}
   \operatorname{span}C(A_J)
   \xrightarrow{\;\pi_R\;}
   \operatorname{span}C(B)
   \longrightarrow0.
\]
Hence the effective dimensions satisfy the exact additive identity
\[
 \dim\operatorname{span}C(A_J)
       =\dim\ker A_S+\dim\operatorname{span}C(B).
\]
All matrix constructions can be performed in polynomial rational bit
time.  The theorem includes $S=\varnothing$, with $V=0$ and a vacuous
positive vector.
\end{theorem}

\begin{proof}
First observe that the columns in $S$ positively generate their entire
linear span: if $A_Su=v$ for signed coefficients $u$, then
$A_S(u+cp)=v$, and $u+cp\geq0$ for all sufficiently large $c$.
Thus $\operatorname{cone}(A_S)=V$.

If $(x,z)\in C(A_J)$, applying $Q$ gives $Bz=0$, so the projection is
well-defined.  Conversely, let $z\geq0$ satisfy $Bz=0$.  Then
$A_Rz\in V$, so choose signed $u$ with $A_Su=-A_Rz$.  For sufficiently
large $c$, $u+cp\geq0$ and $(u+cp,z)\in C(A_J)$.  This proves
surjectivity on the cones.  Every coordinate of $S$ remains feasible
by embedding the original positive vector; surjectivity determines
exactly which coordinates of $R$ are feasible.  This proves the
support identity.

The injection in the displayed sequence is well-defined because
$\ker A_S=\operatorname{span}C(A_S)$ and $C(A_S)$ embeds into
$C(A_J)$.  Surjectivity of the last arrow follows by taking spans of
the already established cone surjectivity.  A vector in the middle
space with zero $R$ coordinates has form $(u,0)$, and its matching
equation says $A_Su=0$.  Conversely every such vector belongs to the
image of the first arrow.  This proves exactness and therefore the
dimension formula.

An exact basis of the annihilator of $\operatorname{im}A_S$ provides
$Q$.  Standard rational elimination and its minor bounds give the
stated construction and bit bounds.
\end{proof}

This quotient construction preserves the linear nonnegative cone.
If $J$ contains mutually incompatible quadrilateral types, the lift by
$cp$ may use incompatible types; the theorem then describes a safe
matching relaxation and does not itself solve the compatibility
problem.  For a compatible fixed sector, this particular obstruction
does not occur, although the separate geometric verification remains.

\section{Sparse families separating raw and effective dimension}

\begin{theorem}[Cyclic obstruction to bounded-row sign propagation]
Fix integers $m\geq3$ and $s\geq1$.  Use nonnegative variables
$p_1,\ldots,p_m,z_1,\ldots,z_m,t_1,\ldots,t_{s-1}$, with cyclic
indices on the $z$ variables, and impose
\[
 p_i+z_i-z_{i+1}=0\qquad(1\leq i\leq m).
\]
In the displayed variable ordering the matrix is already in reduced
row-echelon form.  Its rank is $m$ and its raw matching nullity is
$m+s-1$.  Its nonnegative cone has maximal support consisting of all
$z$ and $t$ coordinates and has effective dimension $s$.

Every nonzero nonnegative linear combination of the displayed rows
uses all $m$ rows.  Here ``nonnegative linear combination'' means a
row combination whose resulting coefficients are all nonnegative;
the combination's multipliers are not assumed nonnegative in advance.
Thus any sign-forcing method examining only spans of at most $h<m$
of these rows makes no first step and therefore no progress.
\end{theorem}

\begin{proof}
The $p$ columns form the identity matrix, proving the rank and RREF
claims.  Sum all equations to obtain $\sum_i p_i=0$.  Nonnegativity
forces $p_i=0$ for every $i$, after which the equations require all
$z_i$ to be equal.  The common value and the $s-1$ free $t$ values
parameterize the cone as $\mathbb R_+^s$.  Setting the common value
and all free values to one proves the support statement.

For a row combination with multipliers $\lambda_i$, the $p_i$
coefficient is $\lambda_i$, so a nonnegative resulting row requires
$\lambda_i\geq0$.  Its $z_i$ coefficient is
$\lambda_i-\lambda_{i-1}$.  All these coefficients are nonnegative
and their sum is zero.  They are therefore all zero, forcing all
$\lambda_i$ to be equal.  A nonzero nonnegative result consequently
has every multiplier strictly positive and uses all $m$ rows.
No bounded-row sign step is initially available, so there are no
earlier forced zeros from which an iterative closure could begin.
\end{proof}

An exact joint certificate is particularly small:
\[
 x_p=0,\quad x_z=\mathbf1,\quad x_t=\mathbf1,
 \qquad \lambda=\mathbf1,
 \qquad y_p=\mathbf1,\quad y_z=y_t=0.
\]
Every row has three nonzero entries, and every matrix entry is in
$\{-1,0,1\}$.  For fixed $s$, the raw dimension tends to infinity
while the complete support reduction exposes a fixed-dimensional
feasible cone.

The obstruction is deliberately relative to the specified row family,
including its RREF presentation.  It is not invariant under arbitrary
change of row basis: a basis containing the sum of all rows would
immediately reveal the forced zeros.  That dependence is exactly why
fixed one-row and two-row closure is insufficient.

\begin{example}[A bounded-coefficient variant with effective dimension zero]
Replace the cyclic equations by
\[
 p_i-z_i+2z_{i+1}=0.
\]
Their sum forces all $p_i,z_i$ to vanish, so the nonnegative cone is
zero despite raw nullity $m$.  A nonnegative row consequence has
$\lambda_i\geq0$ and $\lambda_i\leq2\lambda_{i-1}$.  If one
multiplier vanishes, these cyclic inequalities force every multiplier
to vanish.  Thus every nonzero nonnegative consequence again uses all
$m$ rows.  Appending $s$ unconstrained nonnegative columns gives raw
nullity $m+s$ and effective dimension $s$, for any $s\geq0$.
\end{example}

These are structural examples for abstract rational matching systems.
No realization by triangulated knot exteriors, or by the repository's
actual source-derived sector equations, is asserted.  They establish
an algebraic separation and supply regression cases; a geometric
realization theorem would be additional research.

\section{Conflict branching controlled by intrinsic dimension}

\begin{theorem}[A dimension bound for compatibility search]
Let $A$ and $\ell$ be rational, and let $\mathcal P$ be a set of
forbidden coordinate pairs.  Consider the decision problem
\[
 Ax=0,\quad x\geq0,\quad \ell^Tx>0,
 \qquad x_ix_j=0\quad\text{for every }\{i,j\}\in\mathcal P.
\]
Put $d=\dim\operatorname{span}C(A)$, and let $N$ be the bit length of
the rational input and pair list.  Following complete support
restriction, this problem can be solved using at most
$2^{d+1}-1$ additional linear-feasibility queries, each of polynomial
bit complexity in $N$.  Its total bit running time is
$2^{O(d)}\operatorname{poly}(N)$.
In particular $d=O(\log N)$ gives a polynomial bound, and
$d=O((\log N)^2)$ gives $N^{O(\log N)}$.
\end{theorem}

\begin{proof}
Use homogeneity to replace $\ell^Tx>0$ by $\ell^Tx\geq1$.
At a search node retain the matching equations, nonnegativity and
the coordinate-zero equations selected along its path.  Query
feasibility with $\ell^Tx\geq1$.  An infeasible node is rejected.
If a feasible witness has no forbidden pair positive, it solves the
problem.  Otherwise choose a forbidden pair $\{i,j\}$ whose two
coordinates are positive in this witness and branch on $x_i=0$ and
$x_j=0$.  Every compatible feasible vector belongs to at least one
child, so the branching is complete.

Let $C_v$ be the homogeneous matching cone at a feasible node.
The coordinate functionals $x_i$ and $x_j$ are both nonzero on
$\operatorname{span}C_v$, as witnessed by the positive conflicting
coordinates.  Each child cone is contained in the kernel of one of
these nonzero functionals.  Consequently each child has span
dimension at most $\dim\operatorname{span}C_v-1$; the child is also
a proper coordinate face of $C_v$.  At depth $d$ every remaining
cone is zero, so its normalized LP is infeasible.  A binary tree of
depth at most $d$ has at most $2^{d+1}-1$ nodes.

All node systems consist of the original input and coordinate-zero
conditions, so their input bit lengths remain polynomial in $N$.
Rational linear programming supplies the required polynomial cost per
query.  The displayed complexity consequences follow directly.
\end{proof}

It is enough to compute complete support once at the root: its
restricted kernel has dimension exactly $d$.  Recomputing support at
children can improve practical pruning and dimension estimates but is
unnecessary for the bound.  Similarly, a zero-dimensional homogeneous
cone can be rejected before an LP call, improving constants but not
the theorem's stated bound.

Quadrilateral exclusions are represented by the three forbidden pairs
among the three quadrilateral types of each tetrahedron.  The theorem
is a refinement of familiar LP-guided compatibility branching, not an
assertion that such branching was previously unknown.  Its useful
parameter is the intrinsic matching-cone dimension exposed by the
complete support certificate.

This theorem is an algebraic existence result.  An arbitrary positive
linear objective does not imply an essential normal disc.  Even when
the objective is Euler characteristic in genuine standard normal
coordinates, vertex links, spheres, disconnected components and
boundary-parallel discs require the applicable established geometric
arguments and verification steps.  Nor has a logarithmic-squared
bound on $d$ for arbitrary knot-exterior inputs been proved here.
